% ============================================================
% Chapter 3: Hilbert's Problem 2 — The Consistency of Arithmetic
% ============================================================

\chapter{The Consistency of the Axioms of Arithmetic}

\begin{solvedbox}
\textbf{Resolved (1931).} G\"odel's second incompleteness theorem showed that no sufficiently strong, consistent formal system can prove its own consistency---precisely the kind of internal proof Hilbert had hoped for. The problem is solved, but in a negative sense: Hilbert's original goal turned out to be unachievable.
\end{solvedbox}

% -----------------------------------------------------------
\section{Historical Context: The Foundational Crisis}
% -----------------------------------------------------------

By the end of the nineteenth century, mathematics faced a crisis of confidence. The subject had grown substantially in scope and power, but no one could say with certainty that its foundations were secure.

\subsection{The Dream of Rigor}

Since the time of Euclid, mathematicians had organized their knowledge into systems of \emph{axioms}---self-evident truths from which all theorems follow by logical deduction. Geometry had been axiomatized with great precision by Hilbert himself in his \emph{Grundlagen der Geometrie} (1899). But arithmetic---the simplest and most fundamental branch of mathematics---had never received the same treatment.

What would a complete, rigorous axiomatization of arithmetic look like? Two efforts in particular set the stage.

\subsection{Frege's Logicism}

In 1879, the German logician and philosopher Gottlob Frege published his \emph{Begriffsschrift} (``Concept Script''), a revolutionary system of formal logic. Frege's dream was \emph{logicism}: the idea that all of mathematics could be reduced to pure logic. If successful, this would ground arithmetic---and all of mathematics---in the bedrock of logical reasoning.

Over the next two decades, Frege developed an elaborate logical system. His \emph{Grundgesetze der Arithmetik} (``Basic Laws of Arithmetic'') aimed to derive the natural numbers and their properties from purely logical axioms. Then, in 1901, disaster struck. Bertrand Russell wrote to Frege with a devastating observation:

\begin{quote}
\textit{``Let $R$ be the relation of belonging to the class under consideration. Does the class of all classes that do not belong to themselves belong to itself?''}
\end{quote}

This is the now-famous \textbf{Russell's Paradox}. If such a class exists, it belongs to itself if and only if it does not belong to itself---a contradiction. Russell's paradox showed that Frege's system was inconsistent. Frege himself, upon receiving the letter, wrote: ``A scientist can hardly encounter anything more undesirable than to have the foundation collapse just as the work is finished.''

\subsection{The Peano Axioms}

Independently of Frege, Giuseppe Peano had published in 1889 a set of axioms for the natural numbers, stated in a notation that used logical symbols but did not aim for full logical reduction. The Peano axioms remain the standard starting point for the formal theory of arithmetic:

\begin{enumerate}[label=\textbf{P\arabic*}, leftmargin=2em]
    \item $0$ is a natural number.
    \item For every natural number $n$, $S(n)$ (the \emph{successor} of $n$) is a natural number.
    \item There is no natural number whose successor is $0$.
    \item If $S(m) = S(n)$, then $m = n$.
    \item \emph{The Axiom of Induction:} If $\varphi$ is a property such that $\varphi(0)$ holds, and for every $n$, $\varphi(n)$ implies $\varphi(S(n))$, then $\varphi(n)$ holds for all natural numbers $n$.
\end{enumerate}

These five axioms generate all of arithmetic. Addition and multiplication are defined recursively, and all the familiar properties of arithmetic---commutativity, associativity, the distributive law---are proved from the axioms by induction.

But here is the crucial question: \emph{are these axioms consistent?} Can we be sure that no contradiction follows from them?

\subsection{The Desire for an Internal Proof}

This question was not merely philosophical. If Peano's axioms are inconsistent, then \emph{every} statement is provable from them---including false ones. The entire edifice of number theory would collapse.

Hilbert recognized that the axioms of geometry had been proved consistent relative to the axioms of arithmetic (via coordinate geometry). But this merely pushed the question one level deeper: what certifies the consistency of arithmetic itself?

Hilbert's answer---his \emph{finitist program}---was bold. He proposed that the consistency of arithmetic should be proved by \emph{finitary methods}: simple, concrete, combinatorial reasoning about finite strings of symbols, methods so elementary that their reliability could not be in doubt. Such a proof would be an \emph{internal} proof---a proof conducted entirely within methods more basic than those whose consistency was in question.

\begin{quote}
\textit{``To prove the consistency of arithmetic, we must use only such methods as we are certain are reliable. These finitary methods are so concrete and simple that their correctness can be recognized a priori.''}\par\hfill ---David Hilbert, paraphrased
\end{quote}

This was the heart of Problem 2.

% -----------------------------------------------------------
\section{The Problem}
% -----------------------------------------------------------

In his 1900 lecture, Hilbert stated the second problem as follows:

\begin{problem_statement}
Prove that the axioms of arithmetic are consistent---that is, free from contradiction---by methods that are themselves beyond doubt.
\end{problem_statement}

More precisely, Hilbert asked for a proof of the consistency of the natural numbers that uses only \emph{finitary} means. What counts as ``finitary'' was left somewhat open, but Hilbert had in mind reasoning about finite sequences of symbols, elementary combinatorial arguments, and direct verification of specific configurations---methods that involve no appeal to completed infinities or non-constructive principles.

The idea was to establish a hierarchy of certainty:
\begin{enumerate}
    \item \emph{Finitary} reasoning: the most basic and secure.
    \item \emph{Infinitary} methods (involving quantifiers over infinite sets, induction over arbitrary properties): powerful but less secure.
    \item \emph{Set-theoretic} methods (involving the actual infinite): the most powerful but least secure.
\end{enumerate}

A finitary consistency proof for arithmetic would show that the ``infinitary'' arithmetic is as reliable as the finitary methods used to prove it. This would be a validation of Hilbert's program---a guarantee that mathematics, built step by step from the ground up, stands on solid foundations.

% -----------------------------------------------------------
\section{Key Developments}
% -----------------------------------------------------------

The story of Problem 2 spans three decades and involves some of the significant mathematics of the twentieth century.

\subsection{G\"odel's Incompleteness Theorems (1931)}

In 1931, a twenty-five-year-old Austrian logician named Kurt G\"odel published a short paper that influenced the course of mathematics and philosophy. Its title was ``\"Uber formal unentscheidbare S\"atze der \emph{Principia Mathematica} und verwandter Systeme I'' (``On Formally Undecidable Propositions of \emph{Principia Mathematica} and Related Systems I'').

\paragraph{The First Incompleteness Theorem.}

G\"odel showed:

\begin{theorem}[G\"odel's First Incompleteness Theorem, 1931]
\label{thm:godel-first}
Any consistent, recursively axiomatized formal system $F$ that is capable of expressing basic arithmetic is \emph{incomplete}: there exist statements in the language of $F$ that are neither provable nor refutable in $F$.
\end{theorem}

In other words, no matter how we choose our axioms (as long as they are consistent, can be listed by a computer program, and are strong enough to do arithmetic), there will always be true statements about the natural numbers that the system cannot prove.

The proof is a sophisticated example of self-reference. G\"odel showed how to encode statements \emph{about} the system as statements \emph{within} the system---a technique now called \emph{G\"odel numbering}. He then constructed a sentence $G$ that effectively says: ``This sentence is not provable in $F$.'' If $F$ is consistent, then $G$ is true but unprovable.

\paragraph{The Second Incompleteness Theorem.}

Even more significant for Hilbert's program was G\"odel's second result:

\begin{theorem}[G\"odel's Second Incompleteness Theorem, 1931]
\label{thm:godel-second}
Let $F$ be a consistent, recursively axiomatized system capable of expressing basic arithmetic. Then $F$ cannot prove the statement that expresses its own consistency.
\end{theorem}

The key insight is that the proof of the first theorem can be formalized inside $F$ itself. The system can be shown to be unable to prove its own inconsistency. But by a clever diagonal argument, G\"odel showed that this very fact---that $F$ cannot prove a certain sentence---implies that $F$ cannot prove its own consistency.

Let $\mathrm{Con}(F)$ denote the formalized statement ``$F$ is consistent.'' Then:

\[
    F \text{ is consistent} \;\Longrightarrow\; F \nvdash \mathrm{Con}(F).
\]

This is a precise mathematical theorem, not a philosophical opinion. It shows that any proof of $\mathrm{Con}(F)$ must necessarily use methods \emph{stronger} than $F$ itself---methods that go beyond finitary reasoning.

\begin{remark}
The Peano axioms (P1)--(P5) do not by themselves constitute a recursively axiomatized system, because the induction axiom (P5) is schema---an infinite collection of axioms, one for each formula $\varphi$. However, the first-order theory \emph{Peano Arithmetic} (PA), which takes (P1)--(P4) plus the induction schema, is recursively axiomatized and falls squarely within G\"odel's theorem.
\end{remark}

\subsection{The Significance of G\"odel's Results}

G\"odel's theorems were unexpected. The mathematical community had broadly hoped---as Hilbert had---that a consistency proof could be achieved. The theorems did not say that arithmetic is inconsistent, or that mathematics is unreliable. They said something more subtle: \emph{the certainty of arithmetic cannot be established by arithmetic itself.}

To see why this matters, consider an analogy. Imagine a court system in which each defendant is judged by a jury drawn from the same community. G\"odel's theorem says: if the jury (arithmetic) is consistent, then the jury cannot certify its own consistency. To verify that the jury is fair, you must appeal to a \emph{higher} court.

\subsection{Gentzen's Consistency Proof (1936)}

G\"odel's theorems did not end the story. They shifted it. If a finitary consistency proof for arithmetic is impossible, what happens if we expand the notion of ``acceptable'' methods?

In 1936, Gerhard Gentzen, a student of Hilbert's, accomplished exactly this. He provided a proof of the consistency of Peano Arithmetic---but using a method that went beyond finitary reasoning.

\begin{proposition}[Gentzen, 1936]
\label{prop:gentzen}
Peano Arithmetic is consistent, relative to the assumption of \emph{transfinite induction} up to the ordinal $\varepsilon_0$.
\end{proposition}

What does this mean? Ordinal numbers extend the natural numbers into the transfinite:
\[
    0, 1, 2, \ldots, \omega, \omega+1, \omega+2, \ldots, \omega \cdot 2, \ldots, \omega^2, \ldots, \omega^\omega, \ldots, \omega^{\omega^\omega}, \ldots
\]
The ordinal $\varepsilon_0$ is the smallest ordinal $\alpha$ such that $\omega^\alpha = \alpha$. It is a well-defined, countable ordinal---but it cannot be reached by finitary methods alone.

Gentzen's proof uses \emph{cut elimination} (a technique for simplifying proofs) to reduce any hypothetical proof of $0 = 1$ in Peano Arithmetic to a proof that uses only transfinite induction up to $\varepsilon_0$. Since no such proof exists, arithmetic is consistent---assuming transfinite induction up to $\varepsilon_0$ is reliable.

This is a \emph{relative consistency proof}: it says, ``If you trust transfinite induction up to $\varepsilon_0$, then you may trust Peano Arithmetic.'' It does not give an absolute, self-certifying proof---G\"odel's theorem guarantees that no such proof exists. But it does identify precisely \emph{how much more} than finitary reasoning is needed.

\begin{remark}
Gentzen's result is often summarized as: ``The consistency of PA requires the ordinal $\varepsilon_0$.'' More precisely, the \emph{proof-theoretic ordinal} of Peano Arithmetic is $\varepsilon_0$. This ordinal measures the ``proof-theoretic strength'' of the system and has become a key concept in proof theory.
\end{remark}

% -----------------------------------------------------------
\section{What This Means for Mathematics}
% -----------------------------------------------------------

The resolution of Hilbert's Problem 2 is one of the most significant events in the history of mathematics. Several lessons emerge:

\begin{enumerate}
    \item \textbf{Mathematics cannot fully justify itself from within.} G\"odel's theorem shows that any sufficiently powerful consistent system contains truths it cannot prove, and cannot prove its own consistency. The dream of a self-sufficient, self-certifying foundation for mathematics is impossible.

    \item \textbf{Consistency proofs are possible---at a cost.} Gentzen showed that consistency of arithmetic \emph{can} be proved, but only by using methods that go beyond arithmetic itself. The cost is always an assumption: to prove $T$ consistent, you need a system $T'$ stronger than $T$.

    \item \textbf{The finitist program was partially vindicated.} Hilbert's hope for an absolute, finitary proof was dashed. But the \emph{methodology} of proof theory---formalizing consistency questions, measuring the strength of systems, identifying the precise assumptions needed---became one of the major achievements of twentieth-century mathematics.

    \item \textbf{The foundations of mathematics remain debated.} The impossibility of absolute justification means that mathematicians must, in some sense, \emph{choose} their axioms. This choice is guided by intuition, utility, and fruitfulness---not by a mechanical procedure that can guarantee consistency.
\end{enumerate}

For the undergraduate reader, the essential takeaway is this: the incompleteness theorems are not a statement about the \emph{failure} of mathematics. They are a significant result about its \emph{limits}---limits that, once understood, open the door to a richer understanding of what mathematical truth means.

% -----------------------------------------------------------
\section*{Common Questions}
% -----------------------------------------------------------

\begin{description}[font=\bfseries, leftmargin=2em, style=nextline]

\item[Q: Does G\"odel's theorem mean we can't know anything for sure?]\hfill\\
No. G\"odel's incompleteness theorems are about what can be proved \emph{within} a given formal system, not about what is true or what we can know by other means. We can know that $2+2=4$ perfectly well. The theorems say that if a system is strong enough to do arithmetic, it cannot prove its own consistency and will contain true statements it cannot prove. This is a limit on formal proof, not on knowledge.

\item[Q: Why can't PA prove its own consistency but ZFC can?]\hfill\\
ZFC cannot prove its own consistency either! G\"odel's theorem applies to ZFC just as it applies to PA. But ZFC can prove the consistency of PA because ZFC is stronger than PA---it assumes the existence of infinite sets that PA does not. This is always a \emph{relative} consistency: if you trust ZFC, you can trust PA. But you cannot use ZFC to certify ZFC itself.

\item[Q: What's the difference between ``true'' and ``provable''?]\hfill\\
A statement is \emph{provable} if there exists a finite sequence of logical deductions from axioms that ends with that statement. A statement is \emph{true in the standard model} of arithmetic if it accurately describes the actual natural numbers $\{0,1,2,\dots\}$. G\"odel showed there are true statements about arithmetic that are not provable in PA. Being true is a semantic notion (about what is the case); being provable is a syntactic notion (about what follows from axioms).

\item[Q: How can a statement be ``true but not provable''? Who decides what's true?]\hfill\\
``True'' here means ``true in the standard model $\mathbb{N}$''---the actual natural numbers as we intuitively understand them. G\"odel's sentence $G$ says ``$G$ is not provable in PA.'' If PA were to prove $G$, it would prove a statement asserting its own unprovability, making PA inconsistent. So if PA is consistent, $G$ is not provable in PA---and therefore $G$ is true. Truth is determined by the structure of $\mathbb{N}$, not by our proof systems.

\item[Q: If PA is consistent, why do we need a proof of consistency at all?]\hfill\\
Because consistency is not obvious: PA is a very strong system with infinitely many axioms, and it is easy to accidentally introduce a contradiction when designing such a system. Before Gentzen, no one knew with certainty whether PA was consistent. Moreover, the philosophical program of formalism (Hilbert's program) demanded a finitary proof of consistency as a foundation for all of mathematics. Gentzen's proof showed that consistency \emph{can} be proved, but only by stepping outside PA.

\item[Q: Is there a statement that is ``absolutely'' undecidable---one that no reasonable axiom system can settle?]\hfill\\
Possibly. G\"odel's sentence for PA is undecidable in PA but decidable in ZFC. But there are statements that are independent of ZFC (like CH, discussed in Chapter~2). Whether there are statements that are independent of all ``reasonable'' extensions of ZFC is an open question. Some set theorists believe that CH might be absolutely undecidable, though this is deeply controversial.

\end{description}

% -----------------------------------------------------------
\section{Exercises}
% -----------------------------------------------------------

\begin{enumerate}

    \item \textbf{Consistency.} A formal system is \emph{consistent} if there is no statement $\varphi$ such that both $\varphi$ and $\neg\varphi$ are provable. Explain in your own words why an inconsistent system is useless (i.e., why every statement becomes provable in an inconsistent system). \emph{Hint:} What can you prove if you have both $\varphi$ and $\neg\varphi$?

    \item \textbf{Formal systems.} Consider the following tiny formal system:
    \begin{itemize}
        \item \textbf{Alphabet:} $\{0, S, +, =\}$
        \item \textbf{Axioms:} $0 = 0$ and $S(0) \neq 0$.
        \item \textbf{Rules:} From $a = b$, you may infer $S(a) = S(b)$.
    \end{itemize}
    List all the statements this system can prove in the form $a = b$ where $a$ and $b$ are terms built from $0$ and $S$. Is this system consistent? Is it complete?

    \item \textbf{Self-reference.} G\"odel's proof relies on a sentence that says, in effect, ``I am not provable.'' Explain informally why such a sentence, if it exists in a consistent system, must be true. What would happen if it were false?

    \item \textbf{Gentzen's ordinal.} The ordinal $\varepsilon_0$ is the smallest $\alpha$ satisfying $\omega^\alpha = \alpha$. Write out the first few elements of the sequence $\omega, \omega^\omega, \omega^{\omega^\omega}, \ldots$ and explain why their limit is a natural candidate for a ``fixed point'' of the exponential map $\alpha \mapsto \omega^\alpha$.

    \item \textbf{Philosophical reflection.} G\"odel's theorem says that no consistent system strong enough to do arithmetic can prove its own consistency. Does this mean that we have \emph{no} reason to believe arithmetic is consistent? Or does Gentzen's proof provide genuine reassurance? Write a short paragraph (5--10 sentences) giving your perspective.

    \item \textbf{Further reading ($\star$).} Research the concept of \emph{proof-theoretic ordinal}. Why is $\varepsilon_0$ associated specifically with Peano Arithmetic, and not with, say, Presburger Arithmetic (which is complete and consistent)? What is the proof-theoretic ordinal of Presburger Arithmetic? \emph{Hint:} Presburger Arithmetic cannot express multiplication.

\end{enumerate}
